\section{学风建设}

此处用一句引言开篇，概述班级学风建设的整体理念。

此处介绍班级学风建设的整体成果，如获得的学风荣誉等。

\subsection{奖学金}

此处概述班级同学获得各类奖学金的整体情况，包括获奖人数和奖项分布。

\begin{table}[H]
    \centering
    \zihao{-3}
    \begin{tabular}{||c|c||}
        \hline
        姓名   & 奖项             \\
        \hline
        XXX & 国家奖学金       \\
        XXX & 综合优秀奖       \\
        XXX & 综合优秀奖、XX 奖 \\
        XXX & 学业优秀奖       \\
        XXX & 学业进步奖       \\
        XXX & 社会工作奖       \\
        \hline
    \end{tabular}
    \caption{XX 班同学奖学金情况}
\end{table}

\subsection{科研科创}

此处介绍班级同学参与科研项目和科创赛事的情况，包括参与 SRT 项目或学术计划的人数。

\begin{table}[H]
    \centering
    \zihao{-3}
    \begin{tabular}{||c|c||}
        \hline
        \hline
        姓名   & 期刊会议     \\
        \hline
        XXX & 会议 A 在投 \\
            & 会议 B      \\
        \hline
        XXX & 会议 C 在投 \\
        \hline
        \hline
    \end{tabular}
    \caption{XX 班同学科研情况}
\end{table}

\begin{table}[H]
    \centering
    \zihao{-3}
    \begin{tabular}{||c|c||}
    \hline
    \hline
    姓名   & 比赛          \\
    \hline
    XXX & XX 竞赛 \\
    & 总决赛冠军    \\
    \hline
    XXX & XX 竞赛        \\
    & 团体总分第一  \\
    \hline
    \hline
    \end{tabular}
    \caption{XX 班同学科创情况}
\end{table}

\subsection{学习交流}

此处介绍班级内部学习交流的举措，如学习群的建立与日常活跃情况。

此处描述同学们在群内的交流方式和学习氛围。

% \begin{figure}[H]
%     \centering
%     \includegraphics[width=0.45\textwidth]{assets/example.jpg}
%     \includegraphics[width=0.45\textwidth]{assets/example.jpg}
%     \caption{学习群交流截图}
% \end{figure}

\subsection{学科竞赛}

此处介绍班级同学参与学科竞赛并获奖的情况，包括竞赛名称和获奖级别。

\subsection{总结}

此处总结学风建设的整体成果，可以从班级共同努力的精神、荣誉背后的意义等方面展开。

此处升华主题，强调班级团结一心、携手并进的力量。
